\documentclass{article}
\usepackage{ctex}
\usepackage{amsmath}
\usepackage{tikz}
\usepackage{unicode-math}
\usetikzlibrary{arrows.meta, decorations.markings,patterns}
\begin{document}

2.1-5 %三平行金属板$ A, B$和 $C$,面积都是 $200 cm^2,A,B$ 板相距$4.0$ mm,$A,C$ 板相距 $2.0$ mm,$B,C$ 两板都接地.
%如果使$A$板带正电$ 3.0\times10^{-7}$ C,在略去边缘效应时,问 $B$ 板和 $C$ 板上感应电荷各是多少?以地的电势为零,问 $A $板的电势是多少?


$$ S = 200\ \rm{cm}^2 = 200 \times 10^{-4} \rm{m}^2 = 0.02\rm{m}^2 $$

$ A $ 带正电荷 $ Q_A = 3.0 \times 10^{-7} \rm{C} $

$ A $ 与 $ B $ 之间、$ A $ 与 $ C $ 之间各形成一个电容器

因为 $ B $ 和 $ C $ 接地,$U_B=U_C=0$

$$U_A = U_{AB} = U_{AC}$$

平行板电容公式：
$$C = \frac{\varepsilon_0 S}{d},\varepsilon_0 = 8.85 \times 10^{-12} \ \rm{F/m}$$
$$C_{AB} = \frac{\varepsilon_0 S}{d_{AB}} = \frac{8.85\times 10^{-12} \times 0.02}{0.004}= 4.425\times 10^{-11} \ \rm{F}$$
$$C_{AC} = \frac{\varepsilon_0 S}{d_{AC}} = \frac{8.85\times 10^{-12} \times 0.02}{0.002}= 8.85\times 10^{-11} \ \rm{F}$$


$ A $ 板总电荷 $ Q_A $ 分布在两个表面上：  

面对 $ B $ 的表面电荷为 $ Q_{AB} $,面对 $ C $ 的表面电荷为 $ Q_{AC} $,且：
$$Q_{AB} + Q_{AC} = Q_A ,Q_{AB} = C_{AB} U_A, Q_{AC} = C_{AC} U_A$$

$$(C_{AB} + C_{AC}) U_A = Q_A$$
$$U_A = \frac{Q_A}{C_{AB} + C_{AC}}$$


$$C_{AB} + C_{AC} = 1.3275\times 10^{-10} \ \rm{F},U_A = \frac{3.0\times 10^{-7}}{1.3275\times 10^{-10}} \approx 2.26\times 10^{3} \rm{V}$$

$ B $ 板上的感应电荷 $ Q_B $
$$Q_B = - Q_{AB} = - C_{AB} U_A= - C_{AB}\times \frac{Q_A}{C_{AB} + C_{AC}}  =- 1.0\times 10^{-7} \ \rm{C}$$

同理,$ C $ 板感应电荷 $ Q_C $：
$$Q_C = - Q_{AC} = - C_{AC} U_A = -2.0\times 10^{-7} \ \rm{C}$$

$$\boxed{Q_B = -1.0\times 10^{-7}\ \rm{C},\quad Q_C = -2.0\times 10^{-7}\ \rm{C},\quad U_A \approx 2.26\times 10^{3}\ \rm{V}}$$

















\newpage

2.1-6 %点电荷 $q$ 处在导体球壳的中心,壳的内外半径分别为 $R_1$ 和 $R_2$. 求场强和电势的分布,并画出 $E-r$ 和 $U-r$ 曲线.

球壳内半径 $ R_1 $，外半径 $ R_2 $。
由于球对称，所有物理量只与 $ r $ 有关。


导体内部（ $ R_1 < r < R_2 $）电场为零

因为导体内部场为零，高斯面在 $ R_1 < r < R_2 $ 内必须电通量为零

所以内表面总电荷 = $-q$

中心点电荷 $ q $ 在内表面会感应出均匀分布的电荷 $ -q $。  

球壳原本总电荷为 0，所以外表面均匀分布总电荷 = $ +q $

$$\rho_{\text{内}}(r=R_1) = -\frac{q}{4\pi R_1^2}, \quad\rho_{\text{外}}(r=R_2) = +\frac{q}{4\pi R_2^2}.$$



用高斯定理，取同心球面为高斯面。

(a) $ 0 < r < R_1 $
高斯面内只有中心点电荷 $ q $：
$$E \cdot 4\pi r^2 = \frac{q}{\varepsilon_0}$$
$$E(r) = \frac{q}{4\pi\varepsilon_0 r^2}$$


(b) $ R_1 < r < R_2 $
导体内部，电场为零：
$$E(r) = 0.$$

(c) $ r > R_2 $
高斯面内的总电荷 =  $ q $。
$$E(r) = \frac{q}{4\pi\varepsilon_0 r^2}.$$



$$
E(r) =
\begin{cases}
\dfrac{q}{4\pi\varepsilon_0 r^2}, & 0 < r < R_1, \\
0, & R_1 < r < R_2, \\
\dfrac{q}{4\pi\varepsilon_0 r^2}, & r > R_2.
\end{cases}
$$



取无穷远为电势零点：$ U(\infty) = 0 $。

(a) $ r \ge R_2 $
$$
U(r) = \int_{r}^{\infty} E(r') \, dr'
= \int_{r}^{\infty} \frac{q}{4\pi\varepsilon_0 r'^2} dr'
= \frac{q}{4\pi\varepsilon_0 r}.
$$

(b) $ R_1 \le r \le R_2 $
导体等势，所以：
$$
U(r) = U(R_2) = \frac{q}{4\pi\varepsilon_0 R_2}.
$$

(c) $ 0 < r < R_1 $
从 $ r $ 到 $ R_1 $ 有电场 $ E = \frac{q}{4\pi\varepsilon_0 r'^2} $

从 $ R_1 $ 到 $ R_2 $ 电场为零（电势不变），所以：

$$
U(r) = \int_{r}^{R_1} \frac{q}{4\pi\varepsilon_0 r'^2} dr' + U(R_1)= \frac{q}{4\pi\varepsilon_0} \left( \frac{1}{r} - \frac{1}{R_1} \right) + \frac{q}{4\pi\varepsilon_0 R_2}.
$$

$$
U(r) = \frac{q}{4\pi\varepsilon_0} \left( \frac{1}{r} - \frac{1}{R_1} + \frac{1}{R_2} \right).
$$



$$
U(r) =
\begin{cases}
\dfrac{q}{4\pi\varepsilon_0} \left( \dfrac{1}{r} - \dfrac{1}{R_1} + \dfrac{1}{R_2} \right), & 0 < r < R_1, \\
\dfrac{q}{4\pi\varepsilon_0 R_2}, & R_1 \le r \le R_2, \\
\dfrac{q}{4\pi\varepsilon_0 r}, & r \ge R_2.
\end{cases}
$$





\vspace{2em}

\begin{tikzpicture}[scale=1.6]
    % 坐标轴
    \draw[->] (-0.5,0) -- (6,0) node[right] {$r$};
    \draw[->] (0,-0.5) -- (0,3) node[above] {$E(r), U(r)$};

    
    % E(r) - 蓝色
    \draw[domain=0.5:1.99, smooth, samples=100, thick, blue] 
        plot (\x, {1/(\x*\x)});
    \draw[domain=2.01:3.99, smooth, samples=2, thick, blue] 
        plot (\x, 0);
    \draw[domain=4.01:5.5, smooth, samples=100, thick, blue] 
        plot (\x, {1/(\x*\x)});
    
    % U(r) - 红色
    \draw[domain=0.5:1.99, smooth, samples=100, thick, red] 
        plot (\x, {1*(1/\x - 1/4)});
    \draw[domain=2:4, smooth, samples=2, thick, red] 
        plot (\x, 1);
    \draw[domain=4:5.5, smooth, samples=100, thick, red] 
        plot (\x, {1/\x});
    
    % 参考线标注
    \draw[dashed, gray] (2,0) -- (2,3) node[below right] at (2,0){$R_1$};
    \draw[dashed, gray] (4,0) -- (4,3) node[below right] at (4,0){$R_2$};

    
    % 图例
    \draw[blue, thick] (4,2.8) -- (4.5,2.8) node[right, black] {$E(r)$};
    \draw[red, thick] (5,2.8) -- (5.6,2.8) node[right, black] {$U(r)$};
    



\end{tikzpicture}





















\newpage
2.1-7 %在上题中,若 $q = 4 \times 10^{-10} \, \rm{C}, R_1 = 2 \rm{cm}, R_2 = 3 \rm{cm}$, 求:  

%(1) 导体球壳的电势；  

%(2) 离球心 $r = 1 \rm{cm}$ 处的电势；  

%(3) 把点电荷移开球心 $1 \rm{cm}$,求导体球壳的电势.


$q = 4\times 10^{-10} \ \rm{C}, \quad R_1 = 2\ \rm{cm} = 0.02\ \rm{m}, \quad R_2 = 3\ \rm{cm} = 0.03\ \rm{m}$


 (1)
上一题已得：对于 $ R_1 \le r \le R_2 $，导体是等势体，电势为
$$U_1 = \frac{q}{4\pi\varepsilon_0 R_2}= 120 \ \rm{V}$$

 (2) 
$$U(r) = \frac{q}{4\pi\varepsilon_0} \left( \frac{1}{r} - \frac{1}{R_1} + \frac{1}{R_2} \right), \quad 0 < r < R_1$$
$$\frac{1}{r} = 100, \quad \frac{1}{R_1} = 50, \quad \frac{1}{R_2}=\frac{100}{3}$$

$$U =300 \ \text{V}$$




 (3) 点电荷移开球心 $ 1\ \text{cm} $，求导体球壳的电势


由于静电屏蔽，壳外的电场完全由外表面的 $+q$ 均匀分布决定

壳外电场 = 位于球心的 $+q$ 产生的电场。  

因此壳的电势（$ r \ge R_2 $）为：
$$U_1 = \frac{q}{4\pi\varepsilon_0 R_2}$$
$$U_1 = 120\ \text{V} \quad $$


$$
(1)\ 120\ \text{V}, \quad (2)\ 300\ \text{V}, \quad (3)\ 120\ \text{V}
$$






\newpage

2.1-8半径为$R_1$的导体球带有电荷 $q$,球外有一个内,外半径为$R_2,R_3$ 的同心导体球壳,壳上带有电荷 $Q$

(1)求两球的电势$U_{1}$ 和$U_2$

(2)两球的电势差$\Delta U;$

(3)以导线把球和壳连接在一起后,$U_1,U_2$和$\Delta U$分别是多少?

(4)在情形(1),(2)中,若外球接地,$U_1,U_2$和$\Delta U$为多少?

(5) 设外球离地面很远,若内球接地,情况如何?




静电平衡时，电荷分布为：

(1) 内球：电荷 $ q $ 均匀分布在 $ r = R_1 $ 表面。  
(2) 外球壳内表面（$ r = R_2 $）：由于内球带 $ q $，静电感应使外球壳内表面带 $ -q $（高斯定理，取 $ R_1 < r < R_2 $ 的高斯面，内部总电荷 $ q + (-q) = 0 $）。  
(3) 外球壳外表面（$ r = R_3 $）：外球壳总电荷 $ Q $，内表面已有 $ -q $，所以外表面带 $ Q + q $。



 (1)：未连接导线，未接地
 电场分布（高斯定理）
$$
E(r) =
\begin{cases}
0, & 0 < r < R_1, \\
\dfrac{q}{4\pi\varepsilon_0 r^2}, & R_1 < r < R_2, \\
0, & R_2 < r < R_3, \\
\dfrac{Q+q}{4\pi\varepsilon_0 r^2}, & r > R_3.
\end{cases}
$$



 电势计算
取无穷远 $ U(\infty) = 0 $。

(a) 外球壳的电势 $ U_2 $
外球壳是等势体，取 $ r = R_3 $ 处计算：
$$
U_2 = \int_{R_3}^{\infty} E(r) \, dr
= \int_{R_3}^{\infty} \frac{Q+q}{4\pi\varepsilon_0 r^2} \, dr
= \frac{Q+q}{4\pi\varepsilon_0 R_3}.
$$

(b) 内球的电势 $ U_1 $
从 $ r = R_1 $ 到 $ R_2 $ 有电场 $ \frac{q}{4\pi\varepsilon_0 r^2} $，从 $ R_2 $ 到 $ R_3 $ 电场为 0，从 $ R_3 $ 到无穷远有 $ \frac{Q+q}{4\pi\varepsilon_0 r^2} $：
$$
U_1 = \int_{R_1}^{R_2} \frac{q}{4\pi\varepsilon_0 r^2} \, dr
+ \int_{R_2}^{R_3} 0 \, dr
+ \int_{R_3}^{\infty} \frac{Q+q}{4\pi\varepsilon_0 r^2} \, dr.
$$
$$
= \frac{q}{4\pi\varepsilon_0} \left( \frac{1}{R_1} - \frac{1}{R_2} \right) + 0 + \frac{Q+q}{4\pi\varepsilon_0 R_3}.
$$
$$
U_1 = \frac{1}{4\pi\varepsilon_0} \left[ \frac{q}{R_1} - \frac{q}{R_2} + \frac{Q+q}{R_3} \right].
$$



 答案 (1)
$$
\boxed{U_1 = \frac{1}{4\pi\varepsilon_0} \left( \frac{q}{R_1} - \frac{q}{R_2} + \frac{Q+q}{R_3} \right)}
$$
$$
\boxed{U_2 = \frac{Q+q}{4\pi\varepsilon_0 R_3}}
$$



 3. 情形 (2)：两球电势差
$$
\Delta U = U_1 - U_2
= \frac{1}{4\pi\varepsilon_0} \left[ \frac{q}{R_1} - \frac{q}{R_2} + \frac{Q+q}{R_3} - \frac{Q+q}{R_3} \right]
$$
$$
\Delta U = \frac{q}{4\pi\varepsilon_0} \left( \frac{1}{R_1} - \frac{1}{R_2} \right).
$$

$$
\boxed{\Delta U = \frac{q}{4\pi\varepsilon_0} \left( \frac{1}{R_1} - \frac{1}{R_2} \right)}
$$



 4. 情形 (3)：用导线连接内球与外球壳
连接后，内球与外球壳等势，电荷重新分布。  
设连接后内球电荷变为 $ q' $，外球壳总电荷仍为 $ Q $（系统总电荷 $ q+Q $ 守恒）。  
连接后内球与外球壳等势，且它们之间（$ R_1 < r < R_2 $）电场为零 → 高斯面（$ R_1 < r < R_2 $）内总电荷为零 → 外球壳内表面电荷 = $ -q' $。  
外球壳外表面电荷 = 总电荷 $ Q+q $ 减去内表面电荷 $ -q' $？ 不对，总电荷守恒：  
外球壳总电荷 $ Q $，内表面电荷 $ -q' $，所以外表面电荷 $ Q + q' $。  
系统总电荷：内球 $ q' $ + 外球壳 $ Q $ = $ q' + Q $ = 初始总电荷 $ q+Q $ ⇒ $ q' = q $？ 这不对，因为连接后电荷可以流动，内球电荷可以变，但总电荷守恒：  
内球电荷 $ q_1 $，外球壳总电荷 $ Q $，总电荷 $ q_1 + Q = q + Q $ ⇒ $ q_1 = q $ 吗？ 这样就没变化？ 矛盾。

等一下，仔细想：  
初始：内球 $ q $，外球壳 $ Q $，总电荷 $ q+Q $。  
连接后，内球与外球壳导体连通成为一个导体，电荷全部分布在外表面（$ r = R_3 $），因为静电平衡时导体内部无电荷。  
所以内球表面电荷为 0，外球壳内表面（$ r=R_2 $）电荷也为 0，外表面（$ r=R_3 $）电荷为 $ q+Q $。

因此：
$$
E(r) =
\begin{cases}
0, & 0 < r < R_3, \\
\dfrac{q+Q}{4\pi\varepsilon_0 r^2}, & r > R_3.
\end{cases}
$$
电势：
$$
U_1 = U_2 = \frac{q+Q}{4\pi\varepsilon_0 R_3}.
$$
$$
\Delta U = 0.
$$



 答案 (3)
$$
\boxed{U_1 = U_2 = \frac{q+Q}{4\pi\varepsilon_0 R_3}, \quad \Delta U = 0}
$$



 5. 情形 (4)：外球接地（初始情况 (1) 时接地）
外球接地 ⇒ $ U_2 = 0 $。  
外球壳电势为零，但外球壳外表面电荷会变化，因为接地后电荷可流入大地。  

设接地后外球壳外表面电荷为 $ Q' $，内表面仍为 $ -q $（因为内球仍带 $ q $，且 $ R_1 < r < R_2 $ 的高斯面内总电荷为零）。  
外球壳总电荷 = 内表面 $ -q $ + 外表面 $ Q' $ = 初始外球壳电荷 $ Q $？ 不对，接地后外球壳总电荷可以改变，因为电荷可流入大地。  

更合理：接地后 $ U_2 = 0 $，即 $ r = R_3 $ 处电势为零。  
电势由三部分贡献：  
1. 内球 $ q $（在 $ r=R_1 $）  
2. 外球壳内表面 $ -q $（在 $ r=R_2 $）  
3. 外球壳外表面 $ Q_o $（在 $ r=R_3 $）

在 $ r = R_3 $ 处：
$$
U(R_3) = \frac{1}{4\pi\varepsilon_0} \left[ \frac{q}{R_3} + \frac{-q}{R_3} + \frac{Q_o}{R_3} \right] = \frac{Q_o}{4\pi\varepsilon_0 R_3} = 0
$$
$ Q_o = 0 $。

所以外球壳外表面电荷为 0，内表面仍为 $ -q $，外球壳总电荷 = $ -q $。



 电势计算
内球电势 $ U_1 $：
从 $ R_1 $ 到 $ R_2 $：电场 $ \frac{q}{4\pi\varepsilon_0 r^2} $  
从 $ R_2 $ 到 $ R_3 $：电场为 0（导体内部）  
从 $ R_3 $ 到无穷远：电场为 0（因为外表面电荷 0，且无穷远电势 0）  
所以：
$$
U_1 = \int_{R_1}^{R_2} \frac{q}{4\pi\varepsilon_0 r^2} dr + 0 + 0
= \frac{q}{4\pi\varepsilon_0} \left( \frac{1}{R_1} - \frac{1}{R_2} \right).
$$
$$
U_2 = 0.
$$
$$
\Delta U = U_1 - U_2 = \frac{q}{4\pi\varepsilon_0} \left( \frac{1}{R_1} - \frac{1}{R_2} \right).
$$



 答案 (4)
$$
\boxed{U_1 = \frac{q}{4\pi\varepsilon_0} \left( \frac{1}{R_1} - \frac{1}{R_2} \right), \quad U_2 = 0, \quad \Delta U = \frac{q}{4\pi\varepsilon_0} \left( \frac{1}{R_1} - \frac{1}{R_2} \right)}
$$



 6. 情形 (5)：外球离地很远，内球接地
内球接地 ⇒ $ U_1 = 0 $。  
设内球电荷变为 $ q' $，外球壳总电荷仍为 $ Q $（因为外球壳未接地，电荷守恒？ 不对，内球接地后电荷可流入大地，所以内球电荷可变，外球壳总电荷不变？ 要小心：内球接地后，内球与大地是导体连接，内球电荷可以变化，但外球壳与内球之间没有导线连接，所以外球壳总电荷仍为 $ Q $ 吗？ 是的，因为外球壳孤立（除了与内球电容耦合）。  

设内球电荷 $ q' $，外球壳内表面电荷 $ -q' $，外球壳外表面电荷 $ Q + q' $（因为外球壳总电荷 = 内表面 $ -q' $ + 外表面 $ Q_o $ = $ Q $ ⇒ $ Q_o = Q + q' $）。

内球电势 $ U_1 = 0 $：
$$
U_1 = \frac{1}{4\pi\varepsilon_0} \left[ \frac{q'}{R_1} - \frac{q'}{R_2} + \frac{Q + q'}{R_3} \right] = 0.
$$
$$
\frac{q'}{R_1} - \frac{q'}{R_2} + \frac{Q + q'}{R_3} = 0.
$$
$$
q' \left( \frac{1}{R_1} - \frac{1}{R_2} + \frac{1}{R_3} \right) + \frac{Q}{R_3} = 0.
$$
$$
q' = -\frac{Q/R_3}{\frac{1}{R_1} - \frac{1}{R_2} + \frac{1}{R_3}}.
$$

外球壳电势：
$$
U_2 = \frac{Q + q'}{4\pi\varepsilon_0 R_3}.
$$
电势差：
$$
\Delta U = U_1 - U_2 = -U_2.
$$



 答案 (5)
$$
\boxed{q' = -\frac{Q/R_3}{\frac{1}{R_1} - \frac{1}{R_2} + \frac{1}{R_3}}, \quad U_1 = 0, \quad U_2 = \frac{Q + q'}{4\pi\varepsilon_0 R_3}, \quad \Delta U = -U_2}
$$



最终整理：
1. $ U_1, U_2, \Delta U $ 如上。
2. 连接后等势。
3. 外球接地时外表面电荷为零。
4. 内球接地时内球电荷改变。





好的，我们一步步来解这个同轴圆柱形电容的电势分布问题。  



 1. 已知条件
- 内圆柱体：半径 $ R_1 $，电势 $ U_1 $。
- 外圆柱体：内半径 $ R_2 $，电势 $ U_2 $。
- 两圆柱体很长 → 可视为无限长，忽略边缘效应。
- 圆柱间为真空（或均匀介质，这里按真空处理）。
- 求 $ R_1 < r < R_2 $ 处的电势 $ U(r) $。



 2. 电场分布
由高斯定理，对于长度为 $ L $、半径为 $ r $（$ R_1 < r < R_2 $）的同轴圆柱面，内部电荷为内圆柱单位长度电荷 $ \lambda $（设内圆柱带电线电荷密度为 $ +\lambda $，外圆柱内表面感应 $ -\lambda $）。

电场方向沿径向，大小为：
$$
E(r) = \frac{\lambda}{2\pi\varepsilon_0 r}.
$$



 3. 电势差与 $\lambda$ 的关系
内、外圆柱间电压：
$$
U_1 - U_2 = \int_{R_1}^{R_2} E(r) \, dr
= \int_{R_1}^{R_2} \frac{\lambda}{2\pi\varepsilon_0 r} \, dr
= \frac{\lambda}{2\pi\varepsilon_0} \ln\frac{R_2}{R_1}.
$$
所以：
$$
\lambda = \frac{2\pi\varepsilon_0 (U_1 - U_2)}{\ln(R_2/R_1)}.
$$



 4. 电势分布 $ U(r) $
取从内圆柱（$ r = R_1 $，电势 $ U_1 $）到任意 $ r $ 的积分：
$$
U(r) = U_1 - \int_{R_1}^{r} E(r') \, dr'
= U_1 - \frac{\lambda}{2\pi\varepsilon_0} \int_{R_1}^{r} \frac{dr'}{r'}
$$
$$
= U_1 - \frac{\lambda}{2\pi\varepsilon_0} \ln\frac{r}{R_1}.
$$

代入 $\lambda$：
$$
U(r) = U_1 - \frac{U_1 - U_2}{\ln(R_2/R_1)} \cdot \ln\frac{r}{R_1}.
$$



 5. 整理成常用形式
$$
U(r) = U_1 - (U_1 - U_2) \frac{\ln(r/R_1)}{\ln(R_2/R_1)}.
$$
或者写成：
$$
U(r) = \frac{U_1 \ln(R_2/r) + U_2 \ln(r/R_1)}{\ln(R_2/R_1)}.
$$

验证：  
- $ r = R_1 $ 时，$ U(R_1) = U_1 $。
- $ r = R_2 $ 时，$ U(R_2) = U_2 $。



 6. 最终答案
$$
\boxed{U(r) = U_1 - (U_1 - U_2) \frac{\ln(r/R_1)}{\ln(R_2/R_1)}}
$$
或等价形式：
$$
U(r) = \frac{U_1 \ln\frac{R_2}{r} + U_2 \ln\frac{r}{R_1}}{\ln\frac{R_2}{R_1}}, \quad R_1 < r < R_2.
$$








\newpage

2.1-12 同轴传输线是由两个很长且彼此绝缘的同轴金属直圆柱体构成.
设内圆柱体的电势为$U_{1}$,半径为 $R_1$,外圆柱体的电势为 $U_2$,内半径为 $R_2$,求其间离轴为 $r$ 处$(R_1<r<R_2)$ 的电势.











\newpage

在静电场的作用下,金属导体内部的场强在任何条件下都为零?在很薄的金属薄膜(厚度为几十个甚至几个原子层)内呢?
课本阐述到达静电平衡状态后导体内的场强$E$可以等于零的隐含意义:自由电子数量足够多！查阅材料,根据金属自身的特性(自由电子体密度)来讨论这个问题.




\end{document}
 